% A small XeLaTeX starter adapted from Shakthi Kannan's Tamil/English book
% workflow. Keep the original Org-mode example in example/ for conversion.
\documentclass[11pt,openany]{book}
\usepackage{fontspec}
\usepackage{polyglossia}
\setdefaultlanguage{english}
\setmainfont{DejaVu Serif}
\newfontfamily\tamilfont{DejaVu Sans}
\usepackage{geometry}
\geometry{paperwidth=7in,paperheight=9in,margin=1in,twoside}
\usepackage{fancyhdr}
\usepackage{setspace}
\usepackage{hyperref}
\usepackage{titlesec}
\usepackage{tocloft}
\onehalfspacing
\pagestyle{fancy}
\fancyhead{}
\fancyfoot{}
\fancyhead[LE,RO]{\thepage}
\fancyhead[RE,LO]{\slshape\leftmark}
\title{A Tamil and English Book}
\author{Author Name}
\date{\today}
\begin{document}
\frontmatter
\maketitle
\thispagestyle{empty}
\begin{center}
  \Large\tamilfont தமிழ் புத்தகம்\\[1em]
  \normalsize A bilingual book starter built with XeLaTeX.
\end{center}
\cleardoublepage
\tableofcontents
\mainmatter
\chapter{Getting started}
This starter follows the structure of the original Tamil book workflow:
front matter, a table of contents, and independently editable chapters.
Unicode text can be written directly in the source when a suitable font is
selected with \texttt{fontspec}.
\section{A bilingual paragraph}
English and \tamilfont தமிழ் can appear together in the same document.
Replace this chapter with content generated from the Org-mode files in
\texttt{example/}, or add further chapters with \texttt{\\include}.
\chapter{Next steps}
Tune the page size, headers, fonts, and chapter styling in this file. The
original project’s fuller configuration remains in \texttt{styles/}.
\end{document}
